% !TeX program = xelatex
\documentclass[11pt,twoside,leqno]{article}
\usepackage[paperwidth=512bp,paperheight=798bp,inner=43bp,outer=44bp,
 top=82bp,bottom=49bp,headheight=14pt,headsep=38bp]{geometry}
\usepackage{fontspec}
\usepackage{amsmath}
\usepackage{unicode-math}
\setmainfont{Times New Roman}[
 Path=/mnt/c/Windows/Fonts/,Extension=.ttf,
 UprightFont=*,BoldFont=*bd,ItalicFont=*i,BoldItalicFont=*bi]
\setmathfont{TeX Gyre Termes Math}
\usepackage{amsthm}
\usepackage{microtype}
\usepackage{enumitem}
\usepackage{fancyhdr}
\usepackage{titlesec}
\usepackage{needspace}
\usepackage[hidelinks]{hyperref}
\usepackage{bookmark}
\hypersetup{pdftitle={Chapter 10: Sections 10.14--10.21},
 pdfsubject={Real characters, triangular monomials, and real zeros of Dirichlet L-functions},
 pdfauthor={A newly prepared continuation of Apostol's Chapter 10}}
\setlength{\parindent}{1em}
\setlength{\parskip}{0pt}
\setlength{\emergencystretch}{2em}
\raggedbottom
\setlength{\abovedisplayskip}{8pt plus 2pt minus 2pt}
\setlength{\belowdisplayskip}{8pt plus 2pt minus 2pt}
\titleformat{\section}{\normalfont\fontsize{14}{16}\selectfont\bfseries}
 {\thesection}{0.6em}{}
\titlespacing*{\section}{0pt}{19pt plus 3pt minus 2pt}{7pt}
\renewcommand{\thesection}{10.\arabic{section}}
\setcounter{section}{13}
\setcounter{equation}{25}
\newtheoremstyle{apostol}{9pt}{9pt}{\itshape}{}{\bfseries}{}{0.5em}{}
\theoremstyle{apostol}
\newtheorem{theorem}{Theorem}
\renewcommand{\thetheorem}{10.\arabic{theorem}}
\setcounter{theorem}{15}
\newtheorem{lemma}[theorem]{Lemma}
\newtheorem{corollary}[theorem]{Corollary}
\renewcommand{\proofname}{\normalfont\scshape Proof}
\pagestyle{fancy}
\fancyhf{}
\fancyhead[LE,RO]{\small 10: Primitive roots}
\fancyfoot[LE,RO]{\thepage}
\renewcommand{\headrulewidth}{0pt}
\renewcommand{\sectionmark}[1]{\markright{\thesection: #1}}
\setlist[enumerate]{leftmargin=1.7em,itemsep=4pt,topsep=5pt}
\newcommand{\Z}{\mathbb Z}
\newcommand{\cO}{\mathcal O}
\newcommand{\dd}{\,dt}
\newcommand{\fl}[1]{\left[#1\right]}
\newcommand{\ls}[2]{(#1\mid #2)}

\begin{document}
\setcounter{page}{224}

\section{Real characters and real zeros}
\label{sec:start}

The foregoing sections describe the real characters and the primitive
characters modulo a prime power. We now use these results to study the series
\[
 L(\sigma,\chi)=\sum_{n=1}^{\infty}\frac{\chi(n)}{n^\sigma},
 \qquad \sigma>0,
\]
where \(\chi\) is a real nonprincipal character. By the convergence theorem
(Theorem 6.17), this series converges for every real \(\sigma>0\), and in
Section 6.10 we proved that \(L(1,\chi)\ne0\).

A real zero near 1 is called an exceptional zero, or a Landau--Siegel zero.
Siegel's theorem of 1935 implies that if \(\chi\) is a real primitive character
mod \(q\) and \(L(\beta,\chi)=0\), where \(0<\beta<1\), then, for every
\(\epsilon>0\),
\[
 1-\beta>c(\epsilon)q^{-\epsilon},
\]
where \(c(\epsilon)>0\); his proof gives no method for calculating this
constant, so the estimate is ineffective. Moreover, for fixed \(\epsilon>0\),
\(q^{-\epsilon}\) becomes much smaller than \(1/\log q\) as \(q\) grows.
Thus Siegel's theorem does not give a calculable constant \(c>0\) for which
\(1-\beta\ge c/\log q\). Theorem~\ref{thm:main} gives such a result with
\(c=10^{-6}\). Such a zero affects estimates for primes in arithmetic
progressions, is closely related to class numbers of quadratic fields, and
produces the Deuring--Heilbronn phenomenon. Unlike Siegel's class-number method,
the proof below derives
incompatible upper and lower estimates for the weighted primes with
\(\chi(p)=-1\). Remarkably enough, the first uniform exclusion of
Landau--Siegel zeros was obtained by an AI system \cite{openai}. The argument
presented here follows
the methods of \cite{openai,green,stepanovnote}; related accounts include
\cite{barreto,banks,kerr,tao,vandijk}. The essential ideas of the proof can be
explained by elementary methods accessible to readers familiar with the
material of this book. We prove the following effective estimate.

\begin{theorem}\label{thm:main}
Let \(\chi\) be a real primitive character mod \(q\), where \(q\ge3\).
If \(0<\beta<1\) and \(L(\beta,\chi)=0\), then
\begin{equation}\label{eq:main}
 (1-\beta)\log q\ge10^{-6}.
\end{equation}
\end{theorem}

For \(X\ge2\), let
\begin{equation}\label{eq:S}
 S_-(X)=\sum_{\substack{p\le X\\\chi(p)=-1}}\frac{\log p}{p}.
\end{equation}
The proof of Theorem~\ref{thm:main} is based on two estimates for this sum.
In Sections 10.17 and 10.18 we obtain an upper estimate by constructing a
nonzero integer \(D\) with prescribed prime divisors. In Sections 10.19 and
10.20 a zero \(\beta\) near 1 gives a lower estimate incompatible with the
first. We complete the proof in Section 10.21. All constants are explicit.

We first put a real primitive character into a form suitable for the
congruence argument. Here and below \(\ls d p\) denotes the Legendre symbol
of Chapter 9. The purpose of \(d\) is to turn the condition \(\chi(p)=-1\)
into the elementary assertion that \(d\) is a quadratic nonresidue modulo
\(p\). This will make the ring generated by \(\sqrt d\) useful modulo such
primes.

\begin{theorem}\label{thm:quadratic}
For the character in Theorem~\ref{thm:main} there is a nonzero squarefree
integer \(d\ne1\), with \(|d|\le q\), such that
\[
 \chi(p)=\ls d p\qquad(p\text{ odd prime},\ p\nmid q).
\]
\end{theorem}

For example, the primitive real characters modulo 8 correspond to \(d=2\)
and \(d=-2\). For the quadratic character modulo 5, one may take \(d=5\).

\begin{proof}
We first determine the form of \(q\). Suppose first that \(q=2^\alpha\).
Since \(q\ge3\), we have \(\alpha\ge2\). For \(\alpha=2\), the discussion
preceding Theorem 10.15 shows that the unique nonprincipal character mod \(4\)
is primitive. If \(\alpha\ge3\), Theorem 10.13 shows that a real character
\(\chi_{a,c}\) must have
\[
 c=\frac{\varphi(2^\alpha)}2=2^{\alpha-2}
 \quad\hbox{or}\quad
 c=\frac{\varphi(2^\alpha)}4=2^{\alpha-3},
\]
whereas Theorem 10.15 says that \(\chi_{a,c}\) is primitive if and only if
\(c\) is odd. When \(\alpha=3\), the two values of \(c\) are \(2\) and \(1\),
so the latter gives the primitive real characters mod \(8\). When
\(\alpha\ge4\), both values are even, so no real character mod \(2^\alpha\)
is primitive. Hence \(q=4\) or \(q=8\). If \(q\) is not a power of 2,
its odd part is squarefree and its 2-part is 1, 4, or 8. Here is a hint for
the verification left to the reader in Exercise 20. Factor the character into
its prime-power components as at the end of Section 10.13. For an odd
prime-power component, Theorems 10.12 and 10.14 rule out an exponent at least
2; for the 2-power component, use Theorems 10.13 and 10.15, remembering that
the only character mod \(2\) is induced from mod \(1\). Thus, in all cases,
\[
 q=2^a m,\qquad a=0,2,\text{ or }3,\qquad
 m\text{ odd and squarefree}.
\]
For each odd prime \(r\mid m\), the factorization and primitivity criterion at
the end of Section 10.13 show that the corresponding component modulo \(r\) is
primitive. Theorem 10.12 then identifies it as the quadratic character. Put
\(r^*=(-1)^{(r-1)/2}r\). Quadratic reciprocity and
its supplementary law (Theorems 9.8 and 9.4) give
\[
 \ls{r^*}p=\ls p r\qquad(p\ne r,\ p\text{ odd prime}).
\]
For the factor \(2^a\), the possible primitive real characters, on odd
integers \(n\), are
\[
 (-1)^{(n-1)/2},\qquad
 (-1)^{(n^2-1)/8},\qquad
 (-1)^{(n-1)/2+(n^2-1)/8},
\]
with conductors 4, 8, and 8, respectively. These are the characters
constructed in Section 10.11. By Theorems 9.4 and 9.5 their
values at an odd prime are \(\ls{-4}p,\ls8p,\ls{-8}p\), respectively.
We choose \(D_2=1,-4,8\), or \(-8\), as required, and let
\[
 D_0=D_2\prod_{r\mid m}r^*.
\]
Then \(|D_0|=q\) and \(\chi(p)=\ls{D_0}p\) for every odd prime \(p\nmid q\).
We take \(d=D_0\) if \(a=0\), and \(d=D_0/4\) otherwise. Since 4 is a
square modulo every odd prime, this does not change the Legendre symbol.
This \(d\) is squarefree and \(|d|\le q\). If \(m>1\), an odd prime
divides \(d\); if \(m=1\), its possible values are \(-1,2,-2\).
Hence \(d\ne1\), as required.
\end{proof}

\section{Elementary estimates with constants}
\label{sec:estimates}

We shall use the functions \(\psi(x)\) and \(\vartheta(x)\) of Section
4.2 and the Mangoldt function \(\Lambda(n)\) of Section 2.8.
Theorem 4.10 supplies the estimates
\(\vartheta(x)=O(x)\) and
\(\sum_{p\le x}(\log p)/p=\log x+O(1)\).
To obtain the numerical constant in \eqref{eq:main}, we need explicit
versions of these estimates. The following elementary bounds will suffice.

\begin{theorem}\label{thm:primeestimates}
For every real \(x\ge2\),
\begin{equation}\label{eq:primeestimates}
 \vartheta(x)\le\psi(x)\le4x,\qquad
 \log x-5\le\sum_{p\le x}\frac{\log p}{p}\le\log x+5.
\end{equation}
\end{theorem}

\begin{proof}
Let \(v_p(N)\) denote the exponent of \(p\) in the nonzero integer \(N\).
Legendre's identity (Theorem 3.14) gives
\[
 v_p\binom{2m}{m}
 =\sum_{j\ge1}\left(\fl{2m/p^j}-2\fl{m/p^j}\right)
 \qquad(m\ge1).
\]
Every summand is 0 or 1, and it is 1 if \(m<p^j\le2m\). Therefore
\[
 \psi(2m)-\psi(m)\le\log\binom{2m}{m}\le2m\log2.
\]
The last inequality follows from
\(\binom{2m}{m}\le\sum_{j=0}^{2m}\binom{2m}j=2^{2m}\).
We prove \(\psi(n)\le4n\) by induction. It holds at \(n=1\).
For \(n\ge2\), let \(m=\fl{(n+1)/2}<n\). Then
\[
 \psi(n)\le\psi(2m)\le\psi(m)+2m\log2<6m\le4n.
\]
This proves \(\psi(n)\le4n\) for integral \(n\). If \(x\) is real, apply
this result with \(n=\fl x\).

The factorial identity (Theorem 3.12) gives, for an integer \(n\ge1\),
\[
 \log(n!)=\sum_{j\le n}\Lambda(j)\fl{n/j}.
\]
Since \(0\le n/j-\fl{n/j}<1\), it follows that
\[
 \frac{\log(n!)}n\le\sum_{j\le n}\frac{\Lambda(j)}j
 \le\frac{\log(n!)}n+\frac{\psi(n)}n.
\]
The integral comparison
\(n\log n-n+1\le\log(n!)\le n\log n\) now gives
\[
 \log n-1\le\sum_{j\le n}\frac{\Lambda(j)}j\le\log n+4.
\]
The prime powers of exponent at least 2 contribute at most
\begin{align*}
 \sum_p\sum_{a\ge2}\frac{\log p}{p^a}
 &=\sum_p\frac{\log p}{p(p-1)}
 \le2\sum_{j=2}^{\infty}\frac{\log j}{j^2}\\
 &\le2\left(\frac{\log2}{4}
       +\int_2^\infty\frac{\log t}{t^2}\dd\right)<3.
\end{align*}
The integral comparison is valid since \((\log t)/t^2\) decreases
for \(t\ge2\). Consequently the prime sum at an integer \(n\) lies
between \(\log n-4\) and \(\log n+4\). For real \(x\ge2\), we take
\(n=\fl x\). Then \(x/2\le n\le x\), so
\(\log x-1<\log n\le\log x\). This proves
\eqref{eq:primeestimates}.
\end{proof}

\section{The Frobenius congruence and root counting}
\label{sec:frobenius}

For a prime with \(\chi(p)=-1\), the integer \(d\) is a quadratic
nonresidue modulo \(p\). Consequently the two-coordinate ring used below
becomes a field modulo \(p\), and the Frobenius map sends each element to its
conjugate. A differentiated form of this identity will later force powers of
\(p\) to divide the rows of an integer matrix.

Fix the integer \(d\) of Theorem~\ref{thm:quadratic}, and write
\[
 \cO=\Z[\sqrt d]=\{a+b\sqrt d:a,b\in\Z\}.
\]
The coordinates \(a,b\) are unique. If \(\alpha=a+b\sqrt d\), write
\(\alpha'=a-b\sqrt d\). For negative \(d\), the square root is
taken in the complex numbers. Congruence modulo \(p\cO\) means
congruence of both integer coordinates modulo \(p\).
We use \(\cO\) only as a ring of pairs of integers; no theory of algebraic
integers or ideals is needed. Modulo \(p\cO\), it has the \(p^2\) classes
\(a+b\theta\), where \(a,b\) are integers mod \(p\) and \(\theta^2=d\).

If \(F(z,w)\) is a polynomial, \(F^{(j)}(z,w)\) will denote its \(j\)th
derivative with respect to \(z\), with \(w\) held fixed.

\begin{lemma}\label{lem:residuefield}
If \(p\) is odd and \(\chi(p)=-1\), the residue classes modulo \(p\cO\)
form a field of characteristic \(p\).
\end{lemma}

\begin{proof}
Every residue class has the form \(a+b\theta\), where \(a,b\) are
integers mod \(p\) and \(\theta^2=d\). A nonzero class has inverse
\[
 (a+b\theta)^{-1}=\frac{a-b\theta}{a^2-db^2}.
\]
The denominator does not vanish. For if \(b\ne0\), its vanishing would
give \(d=(a/b)^2\), contrary to \(\ls d p=-1\); and if \(b=0\), then
\(a\ne0\). Hence every nonzero class has an inverse, as asserted.
\end{proof}

\begin{theorem}\label{thm:frobenius}
If \(p\) is odd and \(\chi(p)=-1\), then
\begin{equation}\label{eq:frobenius}
 \alpha^p\equiv\alpha'\pmod{p\cO}\qquad(\alpha\in\cO).
\end{equation}
Moreover, for every \(\alpha\in\cO\), if \(G(z,w)\) has integer coefficients
and \(h,j\) are nonnegative integers, then
\begin{equation}\label{eq:stable}
 \bigl(G(z,w)(z^p-w)^h\bigr)^{(j)}\big|_{(\alpha,\alpha')}
       \in p^h\cO.
\end{equation}
\end{theorem}

\begin{proof}
Since \(\binom pj\) is divisible by \(p\) for \(0<j<p\), the binomial
theorem, Fermat's theorem, and Euler's criterion (Theorems 5.19 and 9.2) give
\[
 (a+b\sqrt d)^p\equiv a+b\,d^{(p-1)/2}\sqrt d
 \equiv a-b\sqrt d\pmod{p\cO}.
\]
This proves \eqref{eq:frobenius}. Also
\[
 (z^p-w)^{(j)}=
 \begin{cases}
 p(p-1)\cdots(p-j+1)z^{p-j},&1\le j\le p,\\
 0,&j>p.
 \end{cases}
\]
At \((\alpha,\alpha')\), every derivative of \(z^p-w\), including the
derivative of order zero, is therefore in \(p\cO\). When we apply the
product rule to \(G(z,w)(z^p-w)^h\), every term contains a derivative
of each of the \(h\) factors \(z^p-w\), and hence contains \(h\) factors
of \(p\). This proves \eqref{eq:stable}.
\end{proof}

\noindent\textit{Note.} The derivatives in \eqref{eq:stable} are ordinary
derivatives; they are not divided by \(j!\). Also, \(j\) need not be less
than \(p\).

\begin{lemma}\label{lem:roots}
Let the coefficients of \(F(z,w)\) belong to a field \(K\) of prime
characteristic \(P\), and let \(1\le m\le P\). Suppose the degree of
\(F\) in \(z\) is less than \(P\), and its degree in \(w\) is less than
\(m\). If
\[
 F^{(j)}(a,a^P)=0\qquad(0\le j<m)
\]
at \(P\) distinct elements \(a\in K\), then \(F=0\).
\end{lemma}

\begin{proof}
Suppose \(F\ne0\), and let \(g(z)=F(z,z^P)\). Distinct monomials
\(z^rw^s\) give distinct exponents \(r+Ps\), since \(0\le r<P\).
Indeed, an equality \(r_1+Ps_1=r_2+Ps_2\) first gives \(r_1=r_2\)
modulo \(P\), and then gives \(s_1=s_2\).
Hence \(g\ne0\), and \(\deg g\le Pm-1\).
The derivative of \(z^P\) is zero in characteristic \(P\), so
\[
 g^{(j)}(z)=F^{(j)}(z,z^P).
\]
If \(g(z)=\sum c_j(z-a)^j\), then \(g^{(j)}(a)=j!c_j\). Since
\(j<m\le P\), none of \(0!,1!,\ldots,(m-1)!\) vanishes in \(K\).
The hypotheses therefore imply that
\((z-a)^m\) divides \(g(z)\). The factor theorem, used in the proof of
Lagrange's theorem (Theorem 5.21), now shows that the \(P\) distinct
values of \(a\) give \(\deg g\ge Pm\), a contradiction.
\end{proof}

\section{A determinant from triangular lattice points}
\label{sec:determinant}

We now prove the arithmetic estimate needed for
Theorem~\ref{thm:main}.

\begin{theorem}\label{thm:primeweight}
For \(X\ge2\) and every integer \(k\ge3\),
\begin{equation}\label{eq:primeweight}
 S_-(X)\le\left(\frac12+\frac7{2k}\right)\log X
               +\log q+\log k+23.
\end{equation}
\end{theorem}

The construction has two distinct uses. The largest prime \(P\) will show
that the rows of a matrix are independent, and hence that one integer minor
\(D\) is nonzero with \(P\nmid D\). Every smaller prime \(p<P\) for which
\(\chi(p)=-1\) will then force a large power of \(p\) to divide the same
integer \(D\). Comparing these divisibilities with an upper bound for
\(|D|\) will give the desired estimate for \(S_-(X)\).

Let \(P\) be the largest odd prime not exceeding \(X\) for which
\(\chi(P)=-1\). Thus the capital letter \(P\) denotes the prime used to
prove that \(D\ne0\), whereas a lower-case \(p\) will later denote one of
the smaller prime divisors forced upon \(D\). The construction is needed
only when \(P\ge5\) and \(k\le P\); throughout this section we assume these
inequalities and let \(n=P-1\).

The weight \(r+ks\) is chosen with the later row operation in mind.
Replacing one factor \(z^p\) by \(w\) changes an exponent pair by
\((-p,1)\). When \(p\ge k\), this cannot increase \(r+ks\), so the
following triangular set remains closed under that operation.

Let \(\mathcal T\) be the set of pairs of nonnegative integers satisfying
\[
 \mathcal T=\{(r,s):r+ks\le n\},
\]
and write
\[
 t=|\mathcal T|,\quad U=\sum_{\mathcal T}r,\quad
 V=\sum_{\mathcal T}s,\quad J=\fl{n/k},\quad m=J+1.
\]
Thus \(t\) is the number of allowed monomials, \(U\) and \(V\) are their
total \(z\)- and \(w\)-degrees, \(J\) is the largest possible \(w\)-degree,
and \(m=J+1\) is the number of derivative orders \(0,\ldots,J\) used at
each point.
Since \(n\ge4\), we have \(U>0\). Every polynomial of the form
\[
 \sum_{(r,s)\in\mathcal T}c_{rs}z^rw^s
\]
has degree less than \(P\) in \(z\) and less than \(m\) in \(w\). We shall
use the inequalities
\begin{equation}\label{eq:counts}
 t\le P^2,\qquad nt\le3U,\qquad kV\le U.
\end{equation}
For a fixed \(s\), the reflection \(r\mapsto n-ks-r\) permutes the
integers in the interval \(0\le r\le n-ks\); hence
\[
 2\sum_{r=0}^{n-ks}r=(n-ks)(n-ks+1).
\]
Summing over \(s\), we obtain
\[
 2U=\sum_{s=0}^J(n-ks)(n-ks+1)=nt-kV,
\]
and therefore \(2U+kV=nt\). To prove \(kV\le U\), we write \(r=ka+b\), where
\(0\le b<k\). For each \(b\), the allowed pairs \((a,s)\) form the
triangle \(a+s\le\fl{(n-b)/k}\), which is unchanged when \(a\) and \(s\)
are interchanged. Hence
\[
 U=k\sum a+\sum b\ge k\sum s=kV,
\]
and consequently \(nt=2U+kV\le3U\). Also
\(\mathcal T\subseteq\{0,\ldots,P-1\}^2\), so \(t\le P^2\).

Let \(N=1+\fl{\sqrt P}\), and select \(P\) distinct pairs from the \(N^2\)
pairs of integers \(0\le u,v<N\). To each chosen pair we associate
\(\alpha=u+v\sqrt d\) and its conjugate \(\alpha'\). Their reductions
modulo \(P\cO\) are distinct, since \(N\le P\), and
Theorem~\ref{thm:frobenius} gives \(\alpha'=\alpha^P\) after reduction.
Moreover,
\[
 |\alpha|,|\alpha'|<H,\qquad H=3\sqrt{qP},
\]
because \(N\le\sqrt P+1\le(3/2)\sqrt P\) and
\(N(1+\sqrt{|d|})\le2N\sqrt q\).

The matrix records the first \(m\) derivatives of every allowed monomial at
the chosen conjugate pairs. More precisely, its rows are indexed by the pairs
\((r,s)\in\mathcal T\). Its columns are indexed by the pairs
\((\alpha,j)\), where \(\alpha\) runs through the \(P\) chosen points and
\(0\le j\le J\). In row \((r,s)\) and column \((\alpha,j)\) we place
\begin{equation}\label{eq:entries}
 (z^rw^s)^{(j)}\big|_{(\alpha,\alpha')},
\end{equation}
Each entry belongs to \(\cO\). Replacing an entry \(a+b\sqrt d\) by its two
integer coordinates \(a,b\), we obtain an integer matrix with \(t\) rows
and \(2Pm\) columns.

\begin{lemma}\label{lem:determinant}
This integer matrix has a square minor \(D\) of order \(t\) such that
\[
 D\ne0\qquad\hbox{and}\qquad P\nmid D.
\]
\end{lemma}

\begin{proof}
Suppose the rows are linearly dependent modulo \(P\). Then there are
coefficients \(c_{rs}\), not all zero, such that the polynomial
\[
 F(z,w)=\sum_{(r,s)\in\mathcal T}c_{rs}z^rw^s,
\]
has its first \(m\) derivative values equal to zero at the \(P\) distinct
pairs \((a,a^P)\) in the field \(\cO/P\cO\) of
Lemma~\ref{lem:residuefield}. Since \(m=\fl{(P-1)/k}+1\le P\),
Lemma~\ref{lem:roots} gives \(F=0\). Hence every \(c_{rs}=0\), a
contradiction. Therefore the rows are independent modulo \(P\), and some
\(t\) columns determine a square minor which is nonzero modulo \(P\). Its
determinant is the required integer \(D\).
\end{proof}

\section{An upper bound for \texorpdfstring{\(S_-(X)\)}{S-(X)}}
\label{sec:arithmetic}

\begin{proof}[\normalfont\scshape Proof of Theorem~\ref{thm:primeweight}]
Let \(P\) have the meaning assigned in the foregoing section. If there is
no such odd prime, or if \(P=3\), then
\[
 S_-(X)\le\frac{\log2}{2}+\frac{\log3}{3}<1,
\]
and the theorem follows. If \(k>P\), all the primes occurring in
\(S_-(X)\) are less than \(k\), and Theorem~\ref{thm:primeestimates}
gives \(S_-(X)\le\log k+5\). We may therefore assume that \(P\ge5\)
and \(k\le P\), and we use the determinant \(D\) of
Lemma~\ref{lem:determinant}.

We now estimate this same integer \(D\) in two ways. Bounds for its entries
give an upper bound for \(\log|D|\). On the other hand, each smaller prime
\(p\) with \(\chi(p)=-1\) divides \(D\) to a known order. Prime
factorization then gives a lower bound for \(\log|D|\).

\medskip
\noindent\textit{Size of the determinant.}
A nonzero entry in row \((r,s)\) of \eqref{eq:entries} is
\((r)_j\alpha^{r-j}(\alpha')^s\), where
\((r)_j=r(r-1)\cdots(r-j+1)\). Since \(r\le n\), \(j\le J\), and
\(H\ge1\), this entry and its conjugate have modulus at most
\(n^JH^{r+s}\). If \(\gamma=a+b\sqrt d\) and \(\gamma'\) both have
modulus at most \(M\), then
\[
 |a|=\left|\frac{\gamma+\gamma'}2\right|\le M,
 \qquad
 |b|=\left|\frac{\gamma-\gamma'}{2\sqrt d}\right|
       \le\frac M{\sqrt{|d|}}\le M.
\]
The same bound therefore holds for both integer coordinates. Index the rows
by \(i=1,\ldots,t\), and let \((r_i,s_i)\) be the pair in row \(i\). If
\(a_{i\ell}\) is any integer entry in that row, the preceding estimate gives
\[
 |a_{i\ell}|\le n^JH^{r_i+s_i}.
\]
For the \(t\) columns defining \(D\), we have
\(\sum_i r_i=U\) and \(\sum_i s_i=V\). The Leibniz formula now gives
\begin{align*}
 |D|
 &=\left|\sum_{\pi\in\mathfrak S_t}\operatorname{sgn}(\pi)
       \prod_{i=1}^t a_{i,\pi(i)}\right|\\
 &\le\sum_{\pi\in\mathfrak S_t}\prod_{i=1}^t|a_{i,\pi(i)}|
 \le t!\prod_{i=1}^t n^JH^{r_i+s_i}
 =t!\,n^{Jt}H^{U+V}.
\end{align*}
Using \(\log(t!)\le t\log t\), \(J\le n/k\), and \eqref{eq:counts},
we obtain
\begin{align}
 \frac{\log|D|}{U}
 &\le\left(1+\frac1k\right)\log H
       +\frac3k\log n+\frac{6\log P}{P-1}\notag\\
 &\le\left(\frac12+\frac7{2k}\right)\log P+\log q+5.
 \label{eq:size}
\end{align}
For the last inequality we used
\(\log H=\tfrac12\log P+\tfrac12\log q+\log3\),
\((1+1/k)\log3<2\), and \(\log P/(P-1)\le1/2\). The coefficient of
\(\log q\) is \(\tfrac12(1+1/k)\le2/3\).

\Needspace{14\baselineskip}
\medskip
\noindent\textit{Divisibility by a smaller bad prime.}
Let \(p\) be a prime such that \(k\le p<P\) and \(\chi(p)=-1\). For
each row we write \(r=ph+c\), where \(h=\fl{r/p}\) and \(0\le c<p\),
and replace its monomial by
\begin{equation}\label{eq:rowchange}
 z^c(z^p-w)^hw^s
 =z^rw^s+\sum_{i=1}^h(-1)^i\binom hi z^{r-pi}w^{s+i}.
\end{equation}
For example, if \(r=7\) and \(p=3\), the row change is
\[
 z^7w^s\longmapsto z(z^3-w)^2w^s
   =z^7w^s-2z^4w^{s+1}+zw^{s+2}.
\]
The two additional monomials have smaller \(z\)-degree; when
\(\chi(3)=-1\), every evaluated coordinate in the new row is divisible by
\(3^2\).
All the new exponent pairs remain in \(\mathcal T\), since
\[
 (r-pi)+k(s+i)=r+ks-i(p-k)\le n,
\]
and they have smaller \(z\)-degree. Hence, when the rows are ordered by
increasing \(r\), the replacements are given by an integer lower triangular
matrix with diagonal entries 1 and leave the fixed minor \(D\) unchanged.
By \eqref{eq:stable}, both integer coordinates in the new row are divisible
by \(p^h\). Factoring these powers from the rows gives
\[
 v_p(D)\ge\sum_{(r,s)\in\mathcal T}\fl{r/p}\ge U/p-t.
\]

\medskip
\noindent\textit{Comparison of the two estimates.}
Since \(D\) is a nonzero integer, unique factorization gives
\[
 \log|D|=\sum_{\ell\mid D}v_\ell(D)\log\ell,
\]
where the sum is over the prime divisors of \(D\). Retaining only the primes
\(k\le p<P\) for which \(\chi(p)=-1\), and then using
Theorem~\ref{thm:primeestimates}, we obtain
\begin{align*}
 \log|D|
 &\ge U\sum_{\substack{k\le p<P\\\chi(p)=-1}}\frac{\log p}{p}
       -t\vartheta(P-1)\\
 &\ge U\sum_{\substack{k\le p<P\\\chi(p)=-1}}\frac{\log p}{p}-12U,
\end{align*}
where we used \((P-1)t\le3U\). The primes below \(k\), including the
prime 2, contribute at most \(\log k+5\), and the prime \(P\) contributes
at most 1. By the choice of \(P\), we have \(S_-(X)=S_-(P)\), and hence
\[
 S_-(X)\le\frac{\log|D|}{U}+12+\log k+5+1.
\]
Together with \eqref{eq:size}, this proves \eqref{eq:primeweight} and
completes the proof.
\end{proof}

\section{A lower bound for \texorpdfstring{\(S_-(X)\)}{S-(X)}}
\label{sec:bridge}

We return to the nonnegative divisor sum of Theorem 6.19. Let
\[
 A(n)=\sum_{a\mid n}\chi(a),\qquad
 B_\beta(x)=\sum_{n\le x}\frac{A(n)}{n^\beta},\qquad 0<\beta<1.
\]
By that theorem, \(A(n)\ge0\) and \(A(1)=1\); consequently
\(B_\beta(x)\ge1\) for \(x\ge1\). We shall use also
\begin{equation}\label{eq:monotonicity}
 A(pn)\ge A(n)\qquad\text{if }\chi(p)\ge0.
\end{equation}
By Theorem 2.14, \(A\) is multiplicative. The prime-power values in the
proof of Theorem 6.19 are \(A(p^e)=e+1\) if
\(\chi(p)=1\) and \(A(p^e)=1\) if \(\chi(p)=0\).
They are nondecreasing in \(e\), which proves
\eqref{eq:monotonicity}, including the case \(p\mid n\).

The following finite-sum inequality relates these auxiliary sums to
the prime weight \(S_-\).

\begin{lemma}\label{lem:bridge}
If \(0<\beta<1\), \(X\ge2\), and \(Y>1\), then
\begin{equation}\label{eq:bridge}
 S_-(X)\ge2\log X
       -\frac{B_\beta(XY)}{B_\beta(Y)}\log(XY)-5.
\end{equation}
\end{lemma}

\begin{proof}
Let \(T=XY\), and denote by \(S_{\ge0}(X)\) the sum
of \((\log p)/p\) over the primes \(p\le X\) with \(\chi(p)\ge0\).
It includes the primes dividing \(q\), and is complementary to
\(S_-(X)\) in the full prime sum. For \(p\le X\),
\(T/p\ge Y\). Since \(A(n)\ge0\), the function \(B_\beta(x)\) is
nondecreasing, and hence \(B_\beta(T/p)\ge B_\beta(Y)\). We also have
\(p^{-\beta}\ge p^{-1}\). Therefore
\begin{align*}
 B_\beta(Y)S_{\ge0}(X)
 &\le\sum_{\substack{p\le X\\\chi(p)\ge0}}
       \frac{\log p}{p^\beta}B_\beta(T/p)\\
 &=\sum_{n\le T}\frac1{n^\beta}
       \sum_{\substack{p\mid n,\ p\le X\\\chi(p)\ge0}}
                    A(n/p)\log p\\
 &\le\sum_{n\le T}\frac{A(n)}{n^\beta}\sum_{p\mid n}\log p
 \le\sum_{n\le T}\frac{A(n)}{n^\beta}\log n.
\end{align*}
Here we used \eqref{eq:monotonicity} and the fact that the product of the
distinct primes dividing \(n\) divides \(n\).

We separate the terms with \(n\le Y\) from those with \(Y<n\le T\).
Using \(\log n\le\log Y\) in the first sum and \(\log n\le\log T\) in the
second, we obtain
\begin{align*}
 \sum_{n\le T}\frac{A(n)}{n^\beta}\log n
 &\le B_\beta(Y)\log Y
   +\bigl(B_\beta(T)-B_\beta(Y)\bigr)\log T\\
 &=B_\beta(T)\log T-B_\beta(Y)\log X.
\end{align*}
This is also an upper bound for \(B_\beta(Y)S_{\ge0}(X)\).
After division by \(B_\beta(Y)>0\), we subtract this inequality from the
lower bound \(\log X-5\) for the full prime sum in
Theorem~\ref{thm:primeestimates}, and obtain \eqref{eq:bridge}.
\end{proof}

\section{Weighted divisor sums at a real zero}
\label{sec:zero}

We next estimate \(B_\beta(x)\) when \(L(\beta,\chi)=0\). As in the proof of
Theorem 6.20, we use the hyperbola identity. We shall need an error estimate
which is uniform for \(3/4\le\beta<1\). The zero will produce a main term
involving \(L(1,\chi)\). The chief point is to combine two remaining tails
before estimating them, so that their factors \(1/(1-\beta)\) cancel.

For \(\sigma>0\), let
\[
 R_\sigma(u)=\sum_{n>u}\frac{\chi(n)}{n^\sigma},\qquad
 L_\sigma(u)=\sum_{n\le u}\frac{\chi(n)}{n^\sigma}.
\]
The periodic cancellation in the proof of Theorem 6.17 gives
\(|C_\chi(t)|\le\varphi(q)\le q\), where
\(C_\chi(t)=\sum_{n\le t}\chi(n)\). If we apply Abel's identity
(Theorem 4.2) and let the upper endpoint tend to infinity, we obtain
\begin{equation}\label{eq:tail}
 R_\sigma(u)=-C_\chi(u)u^{-\sigma}
       +\sigma\int_u^\infty C_\chi(v)v^{-\sigma-1}\,dv,
 \qquad |R_\sigma(u)|\le2qu^{-\sigma}.
\end{equation}
The integral is absolutely convergent. This gives the estimate from
Theorem 6.17 which we shall use.

\Needspace{12\baselineskip}
Let \(\delta=1-\beta\) and
\(H_\beta(v)=\sum_{n\le v}n^{-\beta}\). The proof of Theorem 3.2(b) gives,
more precisely,
\begin{equation}\label{eq:powersum}
 H_\beta(v)=\frac{v^\delta}{\delta}+\zeta(\beta)+e_\beta(v),
 \qquad |e_\beta(v)|\le2v^{-\beta}\quad(v\ge1).
\end{equation}
Here \(\zeta(\beta)\) denotes the constant occurring in Theorem 3.2(b),
not the sum of a convergent infinite series when \(0<\beta<1\).
In fact, if \(\{y\}=y-\fl y\), that proof gives
\[
 e_\beta(v)=-\{v\}v^{-\beta}
   +\beta\int_v^\infty\{y\}y^{-\beta-1}\,dy,
\]
and the estimate follows from \(0\le\{y\}<1\).

\begin{lemma}\label{lem:model}
If \(3/4\le\beta<1\) and \(L(\beta,\chi)=0\), then for \(x\ge1\),
\begin{equation}\label{eq:model}
 B_\beta(x)=\frac{L(1,\chi)}{\delta}x^\delta+E(x),
 \qquad |E(x)|\le10q x^{-1/4}.
\end{equation}
\end{lemma}

\begin{proof}
Let \(u=\sqrt x\). The hyperbola identity (Theorem 3.17), applied to
\(f(a)=\chi(a)a^{-\beta}\) and \(g(b)=b^{-\beta}\), gives
\begin{equation}\label{eq:hyperbola}
 B_\beta(x)=\sum_{a\le u}\frac{\chi(a)}{a^\beta}H_\beta(x/a)
  +\sum_{b\le u}\frac1{b^\beta}L_\beta(x/b)
  -H_\beta(u)L_\beta(u).
\end{equation}
On inserting \eqref{eq:powersum} into the first and last terms, the two
terms \(\zeta(\beta)L_\beta(u)\) cancel. Since \(\beta+\delta=1\),
we obtain
\begin{equation}\label{eq:intermediate}
 B_\beta(x)=\frac{x^\delta}{\delta}L_1(u)
       -\frac{u^\delta}{\delta}L_\beta(u)+E_0(x).
\end{equation}
Here the error term is
\[
\begin{aligned}
 E_0(x)={}&
 \sum_{a\le u}\frac{\chi(a)}{a^\beta}e_\beta(x/a)
 +\sum_{b\le u}\frac1{b^\beta}L_\beta(x/b)\\
 &-e_\beta(u)L_\beta(u).
\end{aligned}
\]
Since \(L(\beta,\chi)=0\), we have \(L_\beta(v)=-R_\beta(v)\), so
\eqref{eq:tail} bounds the three errors by
\begin{align*}
 \left|\sum_{a\le u}\frac{\chi(a)}{a^\beta}e_\beta(x/a)\right|
   &\le2u x^{-\beta}=2u^{1-2\beta},\\
 \left|\sum_{b\le u}\frac1{b^\beta}L_\beta(x/b)\right|
   &\le2qu x^{-\beta}=2qu^{1-2\beta},\\
 |e_\beta(u)L_\beta(u)|&\le4qu^{-2\beta}\le4qu^{1-2\beta}.
\end{align*}
Consequently \(|E_0(x)|\le(2+6q)u^{1-2\beta}\).

The two terms outside \(E_0(x)\) must not be estimated separately: each
contains the large factor \(1/\delta\). Their difference is much smaller.
Indeed, \(R_\beta(u)\) is obtained from the tail \(R_1(u)\) by inserting
the weight \(n^\delta\), and Abel summation exposes a compensating factor
\(\delta\). Put \(a_n=\chi(n)/n\), so that
\(R_1(v)=\sum_{n>v}a_n\). For \(W>u\), Abel's identity (Theorem 4.2),
applied to \(a_n\) and the weight \(v^\delta\), gives
\[
 \sum_{u<n\le W}\frac{\chi(n)}{n^\beta}
 =u^\delta R_1(u)-W^\delta R_1(W)
       +\delta\int_u^W v^{-\beta}R_1(v)\,dv.
\]
By \eqref{eq:tail}, the boundary term at \(W\) tends to zero and
the integral converges absolutely. Therefore
\begin{equation}\label{eq:cancellation}
 \frac{R_\beta(u)-u^\delta R_1(u)}{\delta}
  =\int_u^\infty v^{-\beta}R_1(v)\,dv,
 \qquad
 \left|\int_u^\infty v^{-\beta}R_1(v)\,dv\right|
       \le\frac{2q}{\beta}u^{-\beta}\le3qu^{-\beta}.
\end{equation}
Using \(L_1(u)=L(1,\chi)-R_1(u)\),
\(L_\beta(u)=-R_\beta(u)\), and \(x=u^2\) in
\eqref{eq:intermediate}, we find
\[
 B_\beta(x)=\frac{L(1,\chi)}{\delta}x^\delta
    +u^\delta\int_u^\infty v^{-\beta}R_1(v)\,dv+E_0(x).
\]
The error is at most
\[
 (2+9q)u^{1-2\beta}\le10q x^{1/2-\beta}\le10q x^{-1/4},
\]
because \(q\ge3\) and \(\beta\ge3/4\). This proves
\eqref{eq:model}.
\end{proof}

We have the following consequence.

\begin{theorem}\label{thm:bias}
Under the hypotheses of Lemma~\ref{lem:model}, let \(X\ge2\),
\(Y>1\), and \(\rho=10qY^{-1/4}\). Then
\begin{equation}\label{eq:zerobias}
 S_-(X)\ge2\log X
       -\bigl((1+\rho)X^\delta+\rho\bigr)\log(XY)-5.
\end{equation}
\end{theorem}

\begin{proof}
We apply \eqref{eq:model} at \(XY\) and \(Y\). The coefficient
\(L(1,\chi)/\delta\) cancels, and we obtain
\[
 B_\beta(XY)=X^\delta B_\beta(Y)+E(XY)-X^\delta E(Y).
\]
The two error terms satisfy \(|E(XY)|\le\rho\) and \(|E(Y)|\le\rho\). Since
\(B_\beta(Y)\ge1\), we have
\[
 \frac{B_\beta(XY)}{B_\beta(Y)}
 =X^\delta+\frac{E(XY)-X^\delta E(Y)}{B_\beta(Y)}
 \le X^\delta+\rho+\rho X^\delta
 =(1+\rho)X^\delta+\rho.
\]
Since the coefficient of this ratio in \eqref{eq:bridge} is negative,
the last inequality gives \eqref{eq:zerobias}.
\end{proof}

\section{Exclusion of real zeros near 1}
\label{sec:finish}

We can now complete the proof of Theorem~\ref{thm:main}.
We choose
\begin{equation}\label{eq:parameters}
 Q=10q,\qquad k=64,\qquad X=Q^{128},\qquad Y=Q^8.
\end{equation}
These choices have separate purposes. The value \(Q=10q\) ensures
\(Q\ge30\). The value \(k=64\) makes the leading coefficient in the
determinant bound comfortably smaller than \(2/3\). The choice \(Y=Q^8\)
makes the relative error at most \(1/30\). Finally, \(X=Q^{128}\) is large
enough compared with \(Y\) for the lower estimate to exceed
\((3/4)\log X\), while the contrary assumption still keeps
\((1-\beta)\log X\) small.
Thus \(XY=Q^{136}\), and
\(\rho=10qY^{-1/4}=Q^{-1}\le1/30\).

\begin{proof}[\normalfont\scshape Proof of Theorem~\ref{thm:main}]
Suppose \(L(\beta,\chi)=0\). If \(0<\beta\le3/4\), then
\[
 (1-\beta)\log q\ge\frac14\log3>10^{-6},
\]
and the theorem follows. We may therefore assume that \(3/4<\beta<1\).
Let \(\delta=1-\beta\), and suppose, to the contrary, that
\(\delta\le10^{-6}/\log q\).
Let \(\ell=\log Q>3\). By Theorem~\ref{thm:primeweight},
\begin{align}
 S_-(X)&\le\left(\frac12+\frac7{128}\right)128\ell
                 +\log q+\log64+23\notag\\
       &<72\ell+29<82\ell
        =\frac{41}{64}\log X<\frac23\log X.
        \label{eq:upperfinal}
\end{align}
Here \(\log q<\ell\), \(\log64<6\), and \(29<10\ell\).

On the other hand, \(\log10<3\) and \(\log q>1\) give
\[
 0<\delta\log X
 \le128\cdot10^{-6}\left(1+\frac{\log10}{\log q}\right)
 <512\cdot10^{-6}<\frac1{1000}.
\]
For \(0\le v<1\), comparison of the power series gives
\(e^v\le(1-v)^{-1}\). Thus \(X^\delta\le1000/999\), and
\[
 (1+\rho)X^\delta+\rho
 \le\frac{31}{30}\frac{1000}{999}+\frac1{30}
 =\frac{31999}{29970}<1.07<\frac98.
\]
Theorem~\ref{thm:bias} now gives
\[
 S_-(X)>256\ell-\frac98\,136\ell-5
       =103\ell-5>96\ell=\frac34\log X,
\]
where \(7\ell>5\) was used. This contradicts
\eqref{eq:upperfinal}, and proves the theorem.
\end{proof}

\clearpage
\pdfbookmark[1]{Exercises}{exercises-continued}
\section*{Exercises}
\markright{Exercises}

\small
\begin{enumerate}[start=21,itemsep=7pt,topsep=3pt]
\item Let \(\chi\) be a real nonprincipal character mod \(m\), induced by
a primitive character \(\chi^*\) mod \(q\). Prove that \(L(s,\chi)\) and
\(L(s,\chi^*)\) have the same real zeros for \(0<s<1\). Deduce from
Theorem~\ref{thm:main} that every such zero \(\beta\) satisfies
\[
 (1-\beta)\log m\ge10^{-6}.
\]

\item Under the hypotheses of Lemma~\ref{lem:model}, prove that
\(L(1,\chi)>0\) and that
\[
 B_\beta(x)\sim
 \frac{L(1,\chi)}{1-\beta}\,x^{1-\beta}
 \qquad(x\to\infty).
\]

\item A refinement attributed to K. Barreto improves the dependence on
the conductor in Theorem~\ref{thm:primeweight}. Let \(P\ge5\) be the
largest prime not exceeding \(X\) for which \(\chi(P)=-1\), and suppose
that \(q\le P^2\).
\begin{enumerate}[label=(\alph*),leftmargin=1.8em,itemsep=2pt,topsep=2pt]
\item Choose \(P\) elements \(\alpha=x+y\sqrt d\) of \(\cO\), pairwise
incongruent modulo \(P\cO\), such that
\[
 |\alpha|,|\alpha'|<4P^{1/2}q^{1/4}.
\]
\textit{Hint.} Use a rectangle whose side lengths are close to
\(P^{1/2}q^{1/4}\) and \(P^{1/2}q^{-1/4}\).
\item Use these points in the determinant argument of Sections 10.17--10.18
to prove, for every \(k\ge3\),
\[
 S_-(P)\le\left(\frac12+\frac7{2k}\right)\log P
 +\frac{k+1}{4k}\log q+\log k+O(1),
\]
where the constant in \(O(1)\) is absolute and effective.
\item Take \(k=\lceil\log P\rceil\), treating the bounded values of \(P\)
separately. Split into the cases \(q\le P^2\) and \(q>P^2\), and prove
uniformly for \(X\ge3\) that
\[
 S_-(X)\le\frac12\log X+\frac14\log q+\log\log X+O(1).
\]
\end{enumerate}
\end{enumerate}
\normalsize

\clearpage
\pdfbookmark[1]{Notes and references}{notes-and-references}
\section*{Notes and references}
\markright{Notes and references}

The section numbers, theorem numbers, and equation numbers continue
Chapter 10 of Apostol \cite{apostol}. References to earlier numbered
theorems and exercises are to that book. These sections are newly prepared
text and do not form part of the original 1976 edition.

The terminology and historical summary in Section 10.14 follow the standard
account in Chapter 11 of \cite{montgomery}. Siegel's ineffective estimate
comes from his
1935 class-number paper \cite{siegel}. The repulsion associated with an
exceptional zero grew out of the class-number work of Deuring \cite{deuring}
and Heilbronn \cite{heilbronn}; elementary later treatments include Pintz
\cite{pintz} and Stopple \cite{stopple}.

The nonvanishing principles behind interpolation determinants go back to
Masser \cite{masser} and Laurent \cite{laurent}; the related adelic polynomial
method appears in Bombieri--Vaaler \cite{bombieri}. The multiplicity principle
goes back to Stepanov \cite{stepanov}. The use of an integer minor is attributed
to Barreto \cite{barreto} in the Pescador manuscript. The refinement in
Exercise 23 is likewise attributed there to Barreto; both attributions are
indirect here.

Bostan--Dumas \cite{bostan} discusses Wronskians and linear independence, and
Rosser--Schoenfeld \cite{rosser} gives sharper explicit prime estimates. None
of these deeper results is required above. This presentation claims no new
theorem or priority.

\titlespacing*{\section}{0pt}{8pt}{4pt}
\begin{thebibliography}{99}
\footnotesize\raggedright
\interlinepenalty=10000
\setlength{\itemsep}{0pt}
\setlength{\parskip}{0pt}

\bibitem{openai}
OpenAI (2026). \emph{Uniform exclusion of Landau--Siegel zeros}.
OpenAI Math Release preprint, 1 October 2026.
\href{https://github.com/openai/math/blob/adc7f1241b42e322a6451854ab7e4b4c146bf78a/preprints/Uniform-exclusion-of-Landau-Siegel-zeros-October-1-2026/paper.pdf}
{Manuscript at a fixed revision}.

\bibitem{green}
Green, Ben (2026). On Landau--Siegel zeros. Hexagon preprint, version 2.
\href{https://hexagonmath.org/2610.00311v2}{hexagon:2610.00311v2}.

\bibitem{stepanovnote}
GPT-6 Pro and Claude Opus 5.5 (2026). \emph{Landau--Siegel zeros and a
Stepanov-type nonvanishing lemma}. Preprint posted by J. Pescador, version 2.
\href{https://jacobpescador.github.io/ai-math/papers/landau-siegel-stepanov/}
{Version history and manuscript}.

\bibitem{barreto}
Barreto, K. (2026). \emph{An elementary proof of the non-existence of
Landau--Siegel zeros}. Unpublished manuscript. Cited indirectly through
\cite{stepanovnote} and \cite{vandijk}.

\bibitem{banks}
Banks, William D. (2026). On the OpenAI proof that Landau--Siegel zeros
do not exist. Preprint, version 1.
\href{https://arxiv.org/abs/2610.15483v1}{arXiv:2610.15483v1}.

\bibitem{kerr}
Kerr, Bryce (2026). A simpler proof that there are no Landau--Siegel zeros.
Preprint, version 1.
\href{https://arxiv.org/abs/2610.21099v1}{arXiv:2610.21099v1}.

\bibitem{tao}
Tao, Terence (2026). A sketch of the proof of non-existence of Siegel
zeroes. \emph{What's New}, 19 October 2026.
\href{https://terrytao.wordpress.com/2026/10/19/a-sketch-of-the-proof-of-non-existence-of-siegel-zeroes/}
{Online exposition}.

\bibitem{vandijk}
van Dijk, Rob (2026). A Frobenius determinant proof of the non-existence of
Landau--Siegel zeros. Hexagon preprint, version 2.
\href{https://hexagonmath.org/2610.00409v2}{hexagon:2610.00409v2}.

\bibitem{apostol}
Apostol, Tom M. (1976). \emph{Introduction to Analytic Number Theory}.
New York: Springer-Verlag.

\bibitem{montgomery}
Montgomery, Hugh L., and Robert C. Vaughan (2006).
\emph{Multiplicative Number Theory I. Classical Theory}. Cambridge:
Cambridge University Press.

\bibitem{siegel}
Siegel, Carl Ludwig (1935). \"Uber die Classenzahl quadratischer
Zahlk\"orper. \emph{Acta Arith.} \textbf{1}: 83--86.
\href{https://doi.org/10.4064/aa-1-1-83-86}
{doi:10.4064/aa-1-1-83-86}.

\bibitem{deuring}
Deuring, Max (1933). Imagin\"are quadratische Zahlk\"orper mit der
Klassenzahl 1. \emph{Math. Z.} \textbf{37}: 405--415.
\href{https://doi.org/10.1007/BF01474583}{doi:10.1007/BF01474583}.

\bibitem{heilbronn}
Heilbronn, Hans (1934). On the class-number in imaginary quadratic fields.
\emph{Quart. J. Math. Oxford Ser.} \textbf{5}: 150--160.
\href{https://doi.org/10.1093/qmath/os-5.1.150}
{doi:10.1093/qmath/os-5.1.150}.

\bibitem{pintz}
Pintz, J\'anos (1976). Elementary methods in the theory of \(L\)-functions,
IV: The Heilbronn phenomenon. \emph{Acta Arith.} \textbf{31}: 419--429.
\href{https://eudml.org/doc/205524}{EuDML:205524}.

\bibitem{stopple}
Stopple, Jeffrey (2012). Elementary Deuring--Heilbronn phenomenon.
\emph{Acta Arith.} \textbf{156}: 283--291.
\href{https://doi.org/10.4064/aa156-3-5}{doi:10.4064/aa156-3-5}.

\bibitem{masser}
Masser, David W. (1981). On polynomials and exponential polynomials in
several complex variables. \emph{Invent. Math.} \textbf{63}: 81--95.
\href{https://doi.org/10.1007/BF01389194}{doi:10.1007/BF01389194}.

\bibitem{laurent}
Laurent, Michel (1991). Sur quelques r\'esultats r\'ecents de transcendance.
\emph{Ast\'erisque} \textbf{198--200}: 209--230.
\href{https://www.numdam.org/item/AST_1991__198-199-200__209_0/}
{Numdam full text}.

\bibitem{bombieri}
Bombieri, Enrico, and Jeffrey Vaaler (1983). On Siegel's lemma.
\emph{Invent. Math.} \textbf{73}: 11--32.
\href{https://eudml.org/doc/143033}{EuDML:143033}.

\bibitem{stepanov}
Stepanov, Sergei A. (1969). On the number of points of a hyperelliptic
curve over a finite prime field. \emph{Math. USSR-Izv.} \textbf{3}:
1103--1114.

\bibitem{bostan}
Bostan, Alin, and Philippe Dumas (2010). Wronskians and linear
independence. \emph{Amer. Math. Monthly} \textbf{117}: 722--727.
\href{https://doi.org/10.4169/000298910X515785}
{doi:10.4169/000298910X515785}.

\bibitem{rosser}
Rosser, J. Barkley, and Lowell Schoenfeld (1962). Approximate formulas for
some functions of prime numbers. \emph{Illinois J. Math.} \textbf{6}: 64--94.
\href{https://doi.org/10.1215/ijm/1255631807}
{doi:10.1215/ijm/1255631807}.

\end{thebibliography}

\end{document}
